\documentclass[10pt,a4paper]{amsart}
\usepackage[margin=19mm]{geometry}
\usepackage{amsmath,amssymb,mathtools}
\usepackage[scr=rsfs,cal=euler]{mathalfa}
\usepackage{xfrac}
\usepackage[unicode,hidelinks,bookmarks=false]{hyperref}
\newcommand{\ie}{i.e.\ }
\newcommand{\cf}{cf.\ }
\newcommand{\resp}{respectively\ }
\newcommand{\Subsection}{\subsection}
\providecommand{\nobreakdash}{\nobreak\hbox{-}}
\setlength{\emergencystretch}{2em}
\multlinegap0pt
\usepackage{xcolor}
\usepackage{csquotes}
\newtheorem{thm}{Theorem}
\title{Cycling along Euler road}
\author{Dylan Wyrzykowski}
\date{Source: arXiv:2509.10460v2 (CC BY 4.0)}
\begin{document}
\maketitle

\noindent\emph{Source and licence.} Cycling along Euler road. Version arXiv:2509.10460v2 (CC BY 4.0), distributed by arXiv under the Creative Commons Attribution 4.0 International licence, http://creativecommons.org/licenses/by/4.0/, as recorded in the arXiv licence metadata for this exact version and shown on the abstract page https://arxiv.org/abs/2509.10460v2. Full licence notice: \texttt{licenses/arxiv-2509.10460-CC-BY-4.0.txt}.\par\smallskip

\noindent\textit{Editorial scope.} Counted unit: the article's thm themG (source0.tex lines 159--167). The article's complete original proof of it is delivered with it (source0.tex lines 169--191, one \texttt{proof} environment of 23 lines). No mathematics is written, repaired or re-derived: the statement and the proof are the article's own lines, in the article's own notation, and no retained line is substituted or re-flowed.  No further source-local context is needed: every cross-reference of every retained line resolves inside the two delivered spans.  Nothing was omitted from any retained span, so the package carries no omission record.  No retained line contains a citation, so the wrapper carries no bibliography and no bibliography entry.  No retained line contains a figure environment, an image-inclusion command or drawing code, so no image is delivered and the package declares no asset dependency.  Editorial formatting only, in addition: an AMS \texttt{amsart} preamble that restates the article's own declarations and macros used by the retained lines, namely the environments \texttt{thm} on separate counters, together with the standard packages those lines need.  The retained environments print in this document as Theorem 1 in order of appearance, whereas the article prints the same items with the numbering of its own full document.  The complete native source of the article as distributed in the arXiv e-print for this exact version is bundled unchanged as \texttt{sources/184978/source0.tex}.
\begin{thm}\label{themG}
    Let $A_1, A_2, \ldots, A_n \in \mathbb{S}^{d-1}(O,1)$, with $n > 2, d > 1$. For each $i \in \{1,2,\ldots,n\}$, let $P_{\lambda, i}^{n-1}$ denote the $P_{\lambda}$ point of the sub-polygon $A_1A_2\ldots A_n \setminus A_i$, and let $P_{\lambda}^n$ be the $P_{\lambda}$ point of $A_1A_2\ldots A_n$. Then, the following holds:
    \begin{enumerate}
        \item the multidimensional polygon $P_{\lambda,1}^{n-1}P_{\lambda,2}^{n-1}\ldots P_{\lambda,n}^{n-1}$ is inscribed in the hypersphere $\mathbb{S}^{d-1}(P_{\lambda}^n, |\lambda|)$,
        \item each hypersphere $\mathbb{S}^{d-1}(P_{\lambda, i}^{n-1},|\lambda|)$ passes through $P_{\lambda}^{n}$,
        \item there exists a homothety $\Phi_{\mathbf{T}}^{-\lambda}$ that maps the multidimensional polygon $A_1A_2\ldots A_n$ to the multidimensional polygon $P_{\lambda,1}^{n-1}P_{\lambda,2}^{n-1}\ldots P_{\lambda,n}^{n-1}$,
        \item the lines $A_iP_{\lambda,i}^{n-1}$ concur at a point on the Euler line of $A_1A_2\ldots A_n$.
    \end{enumerate}
\end{thm}

\begin{proof}
    For point 1, we compute the distance between $P_{\lambda}^n$ and each $P_{\lambda,i}^{n-1}$:
\begin{equation}
    \begin{split}
        |p_{\lambda}^n - p_{\lambda,i}^{n-1}| &= \left| \lambda \sum_{j=1}^{n} a_j - \left( \lambda \sum_{j=1}^{n} a_j - \lambda a_i \right)\right| \\
        &= |\lambda a_i| = |\lambda|,
    \end{split}
\end{equation}
    since $A_i \in \mathbb{S}^{d-1}(O,1)$. This proves that all points $P_{\lambda,i}^{n-1}$ lie on the hypersphere $\mathbb{S}^{d-1}(P_{\lambda}^n, |\lambda|)$. For point 2, by symmetry, $P_{\lambda}^n$ lies at distance $|\lambda|$ from each $P_{\lambda,i}^{n-1}$, placing it on each hypersphere $\mathbb{S}^{d-1}(P_{\lambda,i}^{n-1}, |\lambda|)$. For point 3, consider the homothety $\Phi_{\mathbf{T}}^{-\lambda} : \mathbb{R}^d \to \mathbb{R}^d$ defined as follows:
    \begin{equation}\label{homodef}
        \Phi_{\mathbf{T}}^{-\lambda}: x \mapsto -\lambda(x - \mathbf{t}) + \mathbf{t}, \quad \text{where} \quad \mathbf{t} = \frac{\lambda}{\lambda + 1} \sum_{j=1}^{n} a_j.
    \end{equation}
    
    Thus, for any $i \in \{1,2,\ldots,n\}$ we have: \begin{equation}
        \begin{split}
              \Phi^{-\lambda}_{\mathbf{T}}(a_i) &= -\lambda\left(a_i - \frac{\lambda}{\lambda + 1}\sum_{j=1}^{n}a_j\right) + \frac{\lambda}{\lambda + 1}\sum_{j=1}^{n}a_j \\
        &= -\lambda a_i + \frac{\lambda^2}{\lambda + 1}\sum_{j=1}^{n}a_j + \frac{\lambda}{\lambda + 1}\sum_{j=1}^{n}a_j \\
        &= -\lambda a_i + \lambda\left(\frac{\lambda + 1}{\lambda + 1}\sum_{j=1}^{n}a_j\right) \\
        &= \lambda\sum_{j=1}^{n}a_j - \lambda a_i = p_{\lambda,i}^{n-1}.
        \end{split}
    \end{equation}
     Proving point 3. For point 4, since $\Phi_{\mathbf{T}}^{-\lambda}$ maps $A_i$ to $P_{\lambda,i}^{n-1}$, the lines $A_iP_{\lambda,i}^{n-1}$ must all pass through the center of homothety $\mathbf{T}$, which lies on the Euler line of $A_1A_2\ldots A_n$ since $\mathbf{t}= P_{\frac{\lambda}{\lambda +1}}$, by definition \eqref{homodef}.
\end{proof}

\end{document}
